% Draft methods section for the paper, generated 2026-04-21 after the
% Phase-A/B refinements. Reflects the smoothstep pickup ramp, hover-
% equilibrium IC, and fully-local QP structure. Merge into the main
% manuscript as \section{Methods} subsections.

\section{Methods}
\label{sec:methods}

\subsection{Problem Formulation}
A payload of mass $m_L$ is transported along a prescribed trajectory $\mathbf{r}_L^d(t)\in\mathbb{R}^3$ by $N$ quadrotors connected to a common attachment through flexible cables. Let $\mathbf{r}_i\in\mathbb{R}^3$ denote the position of drone $i$, $m_i$ its mass, and $\mathbf{T}_i\in\mathbb{R}^3$ the tension transmitted through cable $i$. The payload translation obeys
\begin{equation}
m_L \ddot{\mathbf{r}}_L = -m_L g\mathbf{e}_z + \sum_{i=1}^{N}\mathbf{T}_i,
\end{equation}
while each drone satisfies
\begin{equation}
m_i \ddot{\mathbf{r}}_i = -m_i g\mathbf{e}_z + f_i R_i\mathbf{e}_z - \mathbf{T}_i,
\end{equation}
with $R_i\in SO(3)$ its attitude and $f_i$ its thrust magnitude. Cable severance is modelled as an instantaneous loss of $\mathbf{T}_k$ for a faulted agent $k\in\mathcal{F}\subset\{1,\dots,N\}$. The control objective is to ensure $\|\mathbf{r}_L(t)-\mathbf{r}_L^d(t)\|$ remains bounded for all admissible fault sets, using only locally-available signals at each drone.

\subsection{Bead-Chain Cable Model}
Each cable is discretised into $N_b=8$ point masses (beads) connected by Kelvin--Voigt links, giving $N_{\mathrm{seg}}=N_b+1=9$ segments. For segment $j$ of rest length $\ell_0=L/N_{\mathrm{seg}}$, the tension-only axial force is
\begin{equation}
F_j = \big[k_s(\|\mathbf{d}_j\|-\ell_0) + c_s\,\hat{\mathbf{d}}_j^\top\dot{\mathbf{d}}_j\big]_+\,\hat{\mathbf{d}}_j,
\end{equation}
where $\mathbf{d}_j$ is the inter-bead vector, $k_s=25\,000$~N/m, $c_s=60$~N$\cdot$s/m, and $[\cdot]_+$ denotes the positive clamp (one-sided rope). The effective end-to-end stiffness is $k_{\mathrm{eff}}=k_s/N_{\mathrm{seg}}\approx 2778$~N/m; the dominant payload-rope mode is
$\omega_n=\sqrt{k_{\mathrm{eff}}/m_L}\approx 30$~rad/s ($\approx 4.8$~Hz) with damping ratio $\zeta\approx 3.7\%$ for the chain-mode seen from the payload. Per-segment dynamics are heavily over-damped ($\zeta_{\mathrm{seg}}\approx 1.2$).

\subsection{Fully-Local Per-Drone Controller}
Each drone solves a three-variable box-constrained quadratic program in its own acceleration command $\mathbf{a}\in\mathbb{R}^3$:
\begin{equation}
\min_{\mathbf{a}}\; w_t\|\mathbf{a}-\mathbf{a}_{\mathrm{tgt}}\|^2 + w_e\|\mathbf{a}\|^2,
\end{equation}
\begin{equation}
\text{s.t.}\;\;|a_x|,|a_y|\le g\tan\theta_{\max},\;\;a_z\in[a_{z,\min},a_{z,\max}].
\end{equation}
The target acceleration embeds a tracking term and a payload-swing damping term,
\begin{equation}
\mathbf{a}_{\mathrm{tgt}} = \mathbf{a}_{\mathrm{trk}} + w_s\,\mathbf{a}_{\mathrm{swing}},\quad \mathbf{a}_{\mathrm{swing}} = -k_s\,\Pi_{xy}\mathbf{v}_L,
\end{equation}
with $\Pi_{xy}=\mathrm{diag}(1,1,0)$ and $\mathbf{v}_L$ the drone-local payload velocity estimate. No peer-state, cable tension of neighbours, or global fault flag enters Eq.~(4)--(6); the controller is thus fully decentralised.

\subsection{Tension Feed-Forward and Pickup Ramp}
The vertical envelope in Eq.~(5) is shaped by the locally-measured cable tension $T_{\mathrm{ff}}$,
\begin{equation}
a_{z,\min}=\tfrac{f_{\min}-T_{\mathrm{ff}}}{m}-g,\qquad a_{z,\max}=\tfrac{f_{\max}-T_{\mathrm{ff}}}{m}-g,
\end{equation}
so that the rope's vertical reaction is cancelled at the feasibility level and a severance on drone $k$ propagates only through $T_{\mathrm{ff},k}\!\to\!0$, making fault response emergent. During pickup, $T_{\mathrm{ff}}$ is gated by a $C^1$ smoothstep,
\begin{equation}
\alpha(\tau)=3\tau^2-2\tau^3,\;\tau=\mathrm{clip}\!\left(\tfrac{t-t_0}{\Delta t},0,1\right),
\end{equation}
satisfying $\alpha'(0)=\alpha'(1)=0$, which replaces the previous linear ramp and eliminates the jerk discontinuity at $\tau\in\{0,1\}$.

\subsection{Hover-Equilibrium Initial Conditions}
To remove startup transients, the simulation is initialised at the analytic fixed point of Eq.~(1)--(3) under zero commanded acceleration. For identical drones lifting a symmetric payload at formation radius $r_f$ with rope rest length $L$, the equilibrium stretch $\delta$ satisfies
\begin{equation}
\delta = \frac{m_L g\,L_{\mathrm{chord}}}{N\,k_{\mathrm{eff}}\,\Delta z_{\mathrm{attach}}},\quad L_{\mathrm{chord}}=L+\delta,\quad \Delta z_{\mathrm{attach}}=\sqrt{L_{\mathrm{chord}}^2 - r_{\mathrm{eff}}^2},
\end{equation}
with $r_{\mathrm{eff}}=0.7\,r_f$ after accounting for the payload-side attachment offset. A fixed-point iteration on Eq.~(9) converges in five iterations to $\delta\approx 3$~mm, $T^\star\approx 8.2$~N per rope, and drone-to-payload vertical offset $1.211$~m.

\subsection{Four-Gate Fault Model}
Cable severance is injected through four logically-separated gates: a \emph{physical} gate that zeroes $\mathbf{T}_k$ in Eq.~(1)--(3); a \emph{telemetry} gate that masks $T_{\mathrm{ff},k}$ in Eq.~(7); a \emph{visual} gate that detaches the rendered cable asset; and a \emph{supervisory} gate that records the event for post-hoc scoring. Each gate is independently toggleable, enabling ablations in which the controller is deliberately denied knowledge that a fault has occurred.

\subsection{Simulation Infrastructure}
The plant is built in Drake as a single \texttt{MultibodyPlant} containing the payload, $N$ bead-chain cables, and $N$ quadrotor bodies imported from URDF and welded to the last bead of their respective cables. Integration uses a semi-implicit Euler scheme at $5$~kHz; the outer controller executes at $5$~kHz and a geometric SO(3) attitude inner loop at $5$~kHz. The 5-drone reference is a three-dimensional Bernoulli lemniscate,
\begin{equation}
x=\tfrac{a\cos\phi}{1+\sin^2\phi},\;y=\tfrac{a\sin\phi\cos\phi}{1+\sin^2\phi},\;z=z_h+A_z\sin\phi,
\end{equation}
with $a=3$~m, $A_z=0.35$~m, and period $T=12$~s.
